\section{Literature Review}
\label{sec:literature-review}

This focused review covers decision levels, exact and approximate methods,
hybridisation and evaluation under uncertainty. It draws on accessible
reviews, methodological papers and retrospective studies; it is not an
exhaustive systematic review. Comparisons concern modelling choices and
evidence rather than a cross-study algorithm ranking.

\subsection{Planning, assignment and sequencing}

Operating room decisions range from strategic capacity to tactical session
allocation and operational case scheduling \citep{cardoen2010,guerriero2011}.
Choosing a date or session does not determine the within-day sequence.
\citet{alamin2025} review deterministic and uncertain formulations, while
\citet{wang2021} distinguish outpatient and inpatient scheduling. That
distinction matters when same-day discharge and overnight stays create
different downstream requirements.

Resource scope changes the problem. \citet{fei2010} separate weekly planning
from daily scheduling with recovery capacity, and \citet{rachuba2026} examine
integrated hospital planning. \citet{carter2024} smooth ward demand through
adjusted surgical activity; \citet{shehadeh2026} use distributionally robust
optimisation for assignment, sequencing and scheduling with uncertain surgery
and downstream stays. These formulations motivate resource integration,
but require information absent from a room-only dataset.

Duration targets must also remain distinct. Room occupancy, surgical
procedure time, turnover and hospital LOS measure different resource uses.
\citet{dupre2025} investigate room occupancy and LOS in emergency surgery;
neither target definition nor predictive performance transfers automatically
to elective planning. Our session-and-bed model and minute-resolution room
model consequently retain separate assumptions and objectives.

\subsection{Exact methods and interpretation of bounds}

Exact methods seek a model-specific optimum and may provide bounds when
interrupted. \citet{cardoen2009} use branch-and-price for day-care sequencing
with recovery resources and multiple criteria. In contrast, \citet{fei2010}
formulate planning as set partitioning but solve it with a column-generation
heuristic before hybrid genetic daily scheduling. An integer formulation
does not by itself certify the resulting solution.

Exhaustive enumeration is useful for small verification instances. With $n$
cases and $m$ rooms, the allocation space contains $m^n$ elements. The gap
$C(s)-C^\star$ is an optimality gap only when $C^\star$ is certified for the
same objective and decoder. An overload-based lower bound can support a bound
gap without proving optimality; a best-known value supports neither claim.
Even an exact predicted optimum does not certify performance under hidden
execution durations.

\subsection{Heuristics and metaheuristics}

Constructive rules provide fast baselines and initial incumbents.
Metaheuristics modify or rebuild these assignments under a finite budget,
with mechanisms that must be interpreted alongside representation and
feasibility handling.

\emph{Simulated annealing} accepts a cost increase $\Delta C>0$ with
probability $\exp(-\Delta C/T)$ as temperature decreases
\citep{kirkpatrick1983}. Temperature is therefore sensitive to objective
scale. \emph{Tabu search} uses move memory and aspiration to guide exploration
beyond local improvement \citep{glover1989}; sampled-neighbour count and
memory length affect both trajectory and effort.

\emph{Genetic algorithms} evolve populations through selection, crossover
and mutation. Their behaviour depends on encoding, repair and incumbent
retention. Surgical examples include \citet{fei2010} and the biased random-key
framework of \citet{vieira2025}. A generation cannot be equated with one
iteration of a single-solution method. \emph{Ant colony optimisation}
constructs solutions using heuristic desirability and pheromone memory
\citep{dorigo1996}. This foundational reference establishes the mechanism,
not hospital-specific superiority. In our allocation adaptation, favouring
a lightly loaded room is only a construction rule; the decoder must enforce
temporal feasibility.

Initial solutions, neighbourhoods, parameter settings and stopping criteria
are needed to explain differences between these methods. Equal fitness-call
budgets control one measure of effort, but candidate construction and
processing costs still differ; runtime must be reported separately.

\subsection{Hybrid approaches and a shared decoder}

Hybridisation can combine search mechanisms or connect distinct modelling
stages. \citet{fei2010} combine planning decomposition and daily search;
\citet{vieira2025} separate a random-key representation from deterministic
schedule decoding. The latter provides a useful comparison for a common
decoder, although our encoding and constraints differ.

The project's Tabu--Simulated Annealing hybrid selects candidates through
tabu memory and applies annealing acceptance. Any benefit requires comparison
with both components under stated budgets and repeated seeds. A lower cost
obtained with more evaluations does not isolate hybridisation itself.

Simulation coupling addresses a different question: whether decisions retain
value when execution diverges from planning assumptions
\citep{guerriero2011,rachuba2026}. Our search evaluates a predicted decoder
cost; Mesa executes the accepted plan and requests replanning at observed
availability events. It does not repeatedly simulate every search candidate.
Reactive policies may consume several request budgets, which must be included
in total effort.

Predictive--prescriptive hybrids add another stage. \citet{ahmadian2025}
combine a neural network with metaheuristics for operating rooms and
post-anaesthesia care under uncertain durations. \citet{kayvanfar2027},
published online in 2026 and assigned to a 2027 volume, connect estimation
with robust multi-resource optimisation. These studies motivate integration
but do not establish that one predictor or solver will dominate in our
simpler models.

\subsection{Prediction, uncertainty and realised performance}

\citet{merghani2025} review machine learning in operating room management,
supporting comparison of nonlinear models with simple historical references.
\citet{dupre2025} report residual uncertainty and speciality differences in
room-occupancy and LOS prediction. Predictive accuracy therefore requires
its own assessment, including temporal validation, calibration and tail errors.

\citet{lex2025} combine duration prediction with elective orthopaedic
optimisation and identify trade-offs between overtime and underutilisation.
\citet{lex2026} study a related single-centre arthroplasty pipeline. These
retrospective designs show how prediction can inform scheduling, while
transfer to other pathways and clinical settings remains conditional.

Prediction and scheduling must respect information timing. Encoders,
imputation and category grouping should use training data; realised durations
remain outcomes, not ordinary solver inputs. An explicit oracle scenario
changes that information boundary and must be labelled separately. A lower
prediction error or predicted objective does not necessarily improve every
operational metric, as the diversity reviewed by \citet{schouten2023} makes
clear. Completion, overtime, delay and resource use remain complementary.

\subsection{Positioning of the present study}

Table~\ref{tab:positioning} summarises the methodological comparisons.
Our contribution is a common experimental interface with explicit limits:
allocation search uses fixed sequencing, session penalties differ from the
room objective, and uncertainty scenarios differ from reactive execution.

\begin{table}[htbp]
\centering\small
\caption{Selected literature and its relevance to the experimental design.
The studies use different resources and are not interchangeable benchmarks.}
\label{tab:positioning}
\begin{tabularx}{\linewidth}{@{}lYY@{}}
\toprule
Study & Relevant approach & Comparison limit \\
\midrule
\citet{cardoen2009} & Exact day-care sequencing. & Richer recovery-resource model. \\
\citet{fei2010} & Weekly planning and hybrid daily scheduling. & Different decomposition and constraints. \\
\citet{vieira2025} & Searches sharing a decoder. & Different representation and search mechanisms. \\
\citet{dupre2025} & Room-occupancy and LOS prediction. & Emergency-surgery cohort and targets. \\
\citet{lex2025,lex2026} & Prediction followed by schedule optimisation. & Single-centre retrospective evidence. \\
\citet{rachuba2026} & Integrated hospital planning. & Broader scope than either project model. \\
\bottomrule
\end{tabularx}
\end{table}
